\documentclass{article}
\usepackage{../../shared/omnicourse}
\input{../../shared/meta.tex}

\renewcommand{\docshortname}{Vision Document}
\hypersetup{pdftitle={OmniCourse -- Vision Document v1.0}}

\begin{document}

\makecover{Vision Document}{\OmniCourseVersion}{\OmniCourseDate}{ESE-PA1-Vision-1.0}{\OmniCourseAssignment}

\pagenumbering{roman}
\thispagestyle{empty}

\section*{Revision history}
\vspace{-0.3em}
{\small
\begin{xltabular}{\dimexpr\textwidth-2\tabcolsep\relax}{@{}p{2.6cm}p{1.8cm}Yp{4.2cm}@{}}
\hcell{Date} & \hcell{Version} & \hcell{Description} & \hcell{Author}\\\hdrline
\endhead
18/Oct/2026 & 1.0 & Initial version submitted for PA1 & Lê Hữu Hoà\\
\end{xltabular}}
\par\vspace{1.0cm}

{\hypersetup{linkcolor=ocNavy}\tableofcontents}

\clearpage
\pagenumbering{arabic}

\section{Introduction}
OmniCourse is a public web application for HCMUS students that puts the degree audit, the curriculum, the semester plan, an
AI planner and the links to the school services on one page. This document defines the high-level needs and the features of
the system: the problem it solves, the stakeholders and the users it serves, the position it takes among the tools that
exist today, the product features and the non-functional requirements. It is written in the Inception phase from the
team's own experience as students, five user interviews and the official curriculum and regulation documents. How each need
is fulfilled in detail belongs to the use-case model and the supplementary specification.

\subsection{References}
\begin{itemize}
  \item Curriculum of the Advanced Program in Computer Science, HCMUS, cohort 2024.
  \item Thông tư 08/2021/TT-BGDĐT of the Ministry of Education and Training on undergraduate training by credits.
  \item Luật 91/2025/QH15 on Personal Data Protection, with Nghị định 356/2025/NĐ-CP.
  \item Course material of Elements of Software Engineering: PA0 project plan, RUP templates, lecture notes on project
        assignments.
  \item Public projects used as prior work: \texttt{LLaammTTeerr/HCMUS\_Courses}, Coursicle, Stellic, Ellucian Degree
        Works, \texttt{nu-map}, \texttt{coursemap}, \texttt{uwpathr}, \texttt{grad-tree}.
\end{itemize}

\section{Positioning}
\subsection{Problem Statement}
The statement below summarises the problem this project solves.

{\small
\begin{xltabular}{\dimexpr\textwidth-2\tabcolsep\relax}{@{}>{\RaggedRight\arraybackslash}p{4.0cm}Y@{}}
\dlabel{The problem of} & academic information that is scattered and unverified. The curriculum, the
requirement groups, the credit limits and the registration steps sit in separate PDF files, separate portals and separate
sets of advice, and each cohort has its own version of them.\\
\dlabel{affects} & HCMUS undergraduates planning their own semesters, in particular the students of a
specialised program such as APCS, and the program office staff who cannot review every student's plan each term.\\
\dlabel{the impact of which is} & a missed compulsory course or an unfilled elective group discovered in
the final year, which delays graduation by a semester, adds tuition and delays a job start; and hours spent each term
re-planning by hand with spreadsheets, screenshots and advice from seniors that nobody has verified.\\
\dlabel{a successful solution would be} & one page that states exactly which requirements are met and
which remain, a roadmap to graduation validated against prerequisites, offering terms and credit caps, the ability to see
the consequence of a choice before making it, and a plain-language explanation of every rule applied.\\
\end{xltabular}}

\subsection{Product Position Statement}
The position statement below summarises the place the product intends to take.

{\small
\begin{xltabular}{\dimexpr\textwidth-2\tabcolsep\relax}{@{}>{\RaggedRight\arraybackslash}p{4.0cm}Y@{}}
\dlabel{For} & HCMUS undergraduates on the credit system, starting with the students of the APCS
program\\
\dlabel{Who} & need to know their real graduation status and plan the semesters ahead without making an
expensive mistake\\
\dlabel{The OmniCourse} & is a web-based academic roadmap and course planner\\
\dlabel{That} & produces a validated plan to graduation and explains every rule it applies, in one
workspace instead of a set of portals, PDF files and spreadsheets\\
\dlabel{Unlike} & hcmus-courses.lamter.cc, the single-purpose registration tools, Coursicle, and the
institutional platforms such as Stellic and Degree Works\\
\dlabel{Our product} & covers the whole cycle in one data model (audit, roadmap, what-if, advisor,
links), keeps every academic rule in our own deterministic engine instead of a language model, and stores curriculum data
in a model that moves between programs and schools.\\
\end{xltabular}}

\section{Stakeholder and User Descriptions}
This section profiles the stakeholders and the users of OmniCourse and the needs they bring to it. Specific requests and
requirements belong to the use-case model and the supplementary specification; this section gives the background and the
justification for them.

\subsection{Stakeholder Summary}
The stakeholders below are not end users of the system, except the development team, which is both.

{\small
\begin{xltabular}{\dimexpr\textwidth-2\tabcolsep\relax}{@{}p{3.9cm}Y Y@{}}
\hcell{Name} & \hcell{Description} & \hcell{Responsibilities}\\\hdrline
\endhead
Instructor and teaching assistants & The staff of the Elements of Software Engineering course &
Review and grade each assignment, approve the project scope, check that the process and the documents follow the course
requirements.\\
Program office and academic affairs staff & The HCMUS staff who own the official curriculum and the academic regulations &
Provide the curriculum and the regulation documents; confirm that the rules applied per cohort are the official ones.\\
Personal data protection officer & The institutional role responsible for student records under Luật 91/2025/QH15 &
Check consent, retention and deletion of personal data, and the notification duty in case of a breach.\\
Development team & Five students of the APCS program, the authors of the system &
Build, test and maintain the product. Also its first users: they plan their own semesters with it and report what does
not work.\\
Students of other universities (later) & Users of the product after this course, not represented in this release &
No responsibility in this project; the data model is designed so that their programs can be added later.\\
\end{xltabular}}

\subsection{User Summary}
Three user groups are identified for this release.

{\small
\begin{xltabular}{\dimexpr\textwidth-2\tabcolsep\relax}{@{}p{3.0cm}Y Y p{3.0cm}@{}}
\hcell{Name} & \hcell{Description} & \hcell{Responsibilities} & \hcell{Stakeholder}\\\hdrline
\endhead
Student (primary) & A HCMUS undergraduate, APCS in the first release, from first year to final year. Plans on a laptop,
checks on a phone. &
Owns their transcript and their plan: imports the transcript, sets the goals, reads the audit, builds a roadmap or asks the
AI planner for one, asks questions in the chat. &
Represented by the development team, who are students of the same program.\\
Program administrator & A member of the program office, or the team during this project, maintaining the curriculum data. &
Adds and edits programs, cohorts, courses, prerequisites, requirement groups and the regulation documents; reviews the data
imported from a curriculum file. &
Represented by the program office stakeholder.\\
Advisor (secondary) & An academic advisor or a lecturer who reviews a student's plan. &
Opens a shared read-only plan and checks it against the curriculum and the regulations of the program. &
Represented by the instructor stakeholder.\\
\end{xltabular}}

\subsection{User Environment}
A student plans on a laptop, in a browser, usually at home in the evening; a quick check happens on a phone between
classes. A planning session runs from ten minutes for a single question to about an hour for a full roadmap to graduation,
and repeats at the start of each term and again around the registration window. One student works alone; a group of friends
may compare plans, which the read-only share covers. The product is a web application, responsive from 360 px to 1920 px,
with a Vietnamese and an English interface, and it needs an internet connection because the AI features call a model API.
It runs in any current browser on Windows, macOS, Android and iOS, with nothing to install.

\begin{itemize}
  \item Systems in use today: the HCMUS portals for registration, grades and tuition, the curriculum PDF, the course
        description files, spreadsheets, and Facebook and Zalo groups.
  \item Integration: none programmatic. OmniCourse links to the portals and never handles a student's password, so nothing
        breaks when the school changes its portal.
  \item Administrators work on a desktop, importing and correcting curriculum data a few times a year, one program at a
        time.
\end{itemize}

\subsection{Summary of Key Stakeholder or User Needs}
The needs below are the ones the product has to satisfy. Priority is set by the cost of getting the need wrong, not by the
ease of building it.

{\small
\begin{xltabular}{\dimexpr\textwidth-2\tabcolsep\relax}{@{}Yp{3.0cm}p{4.0cm}p{1.2cm}@{}}
\hcell{Need} & \hcell{Affects} & \hcell{How it is met today} & \hcell{Priority}\\\hdrline
\endhead
Know the exact graduation status: what is met, what remains, per requirement group & Student &
Reading the transcript by hand, a senior's spreadsheet & High\\
Build a valid plan that respects prerequisites, offering terms and the credit caps & Student &
Spreadsheets, trial and error, asking friends & High\\
Get an answer grounded in the official curriculum and the regulations & Student &
Asking seniors, asking a chatbot with a PDF attached & High\\
Know where and how each school task is performed, and with which account & Student &
Memory, scattered posts in Facebook groups & High\\
See the cost of a choice before making it: fail a course, change specialization, add a summer term & Student &
Advice from peers, guesswork & Medium\\
Maintain curriculum data per program and cohort & Program administrator &
Editing documents, no dedicated tool & Medium\\
Review a student's plan quickly & Advisor & Screenshots and exported spreadsheets & Low\\
\end{xltabular}}

\subsection{Alternatives and Competition}
Students have alternatives today, and most of them are already in use. They are the benchmark this product has to beat.

\begin{itemize}
  \item \textbf{hcmus-courses.lamter.cc} (\texttt{LLaammTTeerr/HCMUS\_Courses}): the closest existing product. Degree plan,
        graduation requirement check and semester planning for HCMUS. \emph{Strength:} it knows the school.
        \emph{Weakness:} it covers one slice of the job and its data model is tied to HCMUS.
  \item \textbf{Automated registration tools} (Auto Course Registration Tool, HCMUS-CTDA-DKHP): they get a student into a
        section faster. \emph{Weakness:} they need the student's portal account, a risk the student carries, and they
        answer none of the planning questions.
  \item \textbf{Coursicle}: schedule building, seat tracking, reviews. \emph{Strength:} polished and proven.
        \emph{Weakness:} one university at a time, no degree audit against a curriculum, no AI planning.
  \item \textbf{Stellic and Ellucian Degree Works}: degree audit and planning platforms. \emph{Strength:} mature audit
        logic. \emph{Weakness:} sold to institutions, not reachable by a student at HCMUS, and closed.
  \item \textbf{Open-source single-school planners} (\texttt{nu-map}, \texttt{coursemap}, \texttt{uwpathr},
        \texttt{grad-tree}): good ideas about roadmap visualisation. \emph{Weakness:} each is one school, no shared data
        model, no AI layer.
  \item \textbf{The status quo}: Excel, screenshots, Facebook and Zalo groups, advice from seniors. \emph{Strength:} free
        and familiar. \emph{Weakness:} unverified, stale, and wrong for the student's own cohort.
\end{itemize}

\section{Product Overview}
\subsection{Product Perspective}
OmniCourse is a new, standalone product, not a module of an existing information system. It is a cloud-hosted web
application: a Vue frontend, a NestJS REST API, a PostgreSQL database, and one language model behind a single interface. It
sits beside the official portals rather than inside them. Integration with the school is by link and by the data the
student provides, which is an uploaded transcript and a chosen program; there is no scraping and no automated transaction
with a student account. The architecture is modular: audit and planning are one module in this release, while tuition,
conduct score, paperwork and community features are independent modules that can be switched on per school later without
changing the engine.

\subsection{Assumptions and Dependencies}
\begin{itemize}
  \item The curriculum and the regulation documents of HCMUS are public and stable within a cohort.
  \item The APCS curriculum can be modelled exactly by the team, who study it.
  \item The credit caps of the program are the published ones: 14 to 25 credits in a main term, 12 in summer.
  \item A student can produce a transcript as a PDF or a screenshot, and consents to upload it.
  \item A hosted model API stays available on a free or low-cost tier for the duration of the course.
  \item Students have a browser and an internet connection; no installation is possible.
  \item No official API for the curriculum or the student record appears during the project.
  \item Accuracy of the curriculum data depends on editorial maintenance by the team after each cohort change.
\end{itemize}

\section{Product Features}
Twelve feature groups deliver the needs above. F1 to F11 serve the student; F12 serves the program administrator. Priority
is set for this course: \emph{High} features must work in the release candidate, \emph{Medium} features are completed if the
earlier sprints hold, \emph{Low} features may be reduced without breaking the product.

{\small
\begin{xltabular}{\dimexpr\textwidth-2\tabcolsep\relax}{@{}p{0.9cm}p{3.9cm}Yp{1.4cm}@{}}
\hcell{No.} & \hcell{Feature} & \hcell{Description} & \hcell{Priority}\\\hdrline
\endhead
F1 & Account and onboarding &
Sign up with email, Google, or a verified student email; choose faculty, program and cohort; set the target graduation
term, the maximum credits per term and a planned internship. & High\\
F2 & Transcript import with AI extraction &
Upload a transcript PDF or image; the model extracts codes, credits, grades and terms; a review screen flags
low-confidence rows; manual entry always possible. & High\\
F3 & Degree audit dashboard &
Progress per requirement group, credits earned and remaining, cumulative GPA, estimated graduation term, and warnings for
an empty elective group, a missing language certificate or a prerequisite not taken. & High\\
F4 & Curriculum explorer &
Every course of the program by requirement group, with credits, prerequisites and what it unlocks, plus a prerequisite
graph. & High\\
F5 & Semester roadmap planner &
Drag courses into terms from now to graduation; every change is validated live against prerequisites, credit caps,
offering terms and requirement coverage; several plans kept and compared. & High\\
F6 & AI plan generator (agentic) &
A plain-language goal becomes a full roadmap: the agent calls the rule engine as tools, repairs violations and explains
each placement. & High\\
F7 & What-if simulator &
Failing a course, changing specialization, adding a summer term or raising the load: the roadmap is recomputed and the
effect on the graduation term is shown. & Medium\\
F8 & AI academic advisor chat &
Questions answered from the official curriculum and regulations, with the document and the clause cited and the student's
own record available. & Medium\\
F9 & Plan export and sharing & Export a plan as a PDF or share a read-only link with an advisor. & Low\\
F10 & Notifications &
Reminders for the registration window, a check-in at the start of each term and a warning when the transcript changes. &
Low\\
F11 & Quicklinks and step-by-step guides &
Every school task the app does not perform comes with its direct link, the steps, the required account and the notes that
matter at that moment; each guide is versioned. & Medium\\
F12 & Curriculum administration &
Create and edit programs, cohorts, courses, prerequisites and requirement groups; import a curriculum PDF into structured
data for review. & Medium\\
\end{xltabular}}

\section{Other Product Requirements}
\subsection{Non-Functional Requirements}
The requirements below apply to the whole product. They are written as targets that can be checked at the end of a sprint,
not as intentions.

\begin{itemize}
  \item \textbf{Security:} passwords hashed with a modern algorithm, optional two-factor authentication, sessions that
        rotate, expire and can be revoked, and administrator rights scoped per school and per program.
  \item \textbf{Privacy:} consent at sign-up, a stated retention period, export and deletion of a whole account, no portal
        credential stored, prompts not used to train third-party models. Basis: Luật 91/2025/QH15 and Nghị định
        356/2025/NĐ-CP.
  \item \textbf{Performance:} the audit and the dashboard load in under 2 seconds for a 40-course transcript on a normal
        connection; a planner edit is validated in under 300 milliseconds; AI plan generation returns in under 30 seconds
        with visible progress.
  \item \textbf{Robustness:} the application stays usable for manual planning when the model provider is unreachable, and
        the database is backed up nightly. Availability target for the course is 99\% monthly on the free tier.
  \item \textbf{Usability:} responsive from 360 px to 1920 px, keyboard operable, Vietnamese and English, colour contrast
        at WCAG 2.1 AA, and every AI answer links to the document and clause it relies on.
  \item \textbf{Testability and maintainability:} at least 80\% statement coverage of the rule engine, automated tests in
        the pipeline on every pull request, and the API contract published as OpenAPI 3.2.0.
  \item \textbf{Data:} curriculum data versioned per cohort, so a new cohort is a new record and old plans keep working.
  \item \textbf{Standards:} national credit regulation Thông tư 08/2021/TT-BGDĐT, the published HCMUS credit caps, and
        Luật 91/2025/QH15 for personal data.
\end{itemize}

\end{document}
